% Author: Grayson Orr
% Course: ID608001: Intermediate Application Development Concepts

\documentclass{article}
\author{}
 
\usepackage{fontspec}
\setmainfont{Arial}

\usepackage{graphicx}
\usepackage{wrapfig}
\usepackage{enumerate}
\usepackage{hyperref}
\usepackage[margin = 1.75cm]{geometry}
\usepackage[table]{xcolor}
\usepackage{fancyhdr}
\hypersetup{
  colorlinks = true,
  urlcolor = blue 
} 
\setlength\parindent{0pt}
\pagestyle{fancy}
\fancyhf{}
\rhead{College of Engineering, Construction and Living Sciences\\Bachelor of Information Technology}
\lfoot{Project \\Version 2, Semester One, 2024}
\rfoot{\thepage}
 
\begin{document}

\begin{figure}
	\centering
	\includegraphics[width=50mm]{../../resources/img/logo.png}
\end{figure}

\title{College of Engineering, Construction and Living Sciences\\Bachelor of Information Technology\\ID608001: Intermediate Application Development Concepts\\Level 6, Credits 15\\\textbf{Project}}
\date{}
\maketitle 

\section*{Assessment Overview}
In this \textbf{individual} assessment, you will develop a three frontend applications. In addition, marks will be allocated for code quality and best practices, documentation and \textbf{Git} usage.

\section*{Learning Outcomes}
At the successful completion of this course, learners will be able to:
\begin{enumerate}
	\item Apply design patterns and programming principles using software development best practices.
	\item Design and implement full-stack applications using industry relevant programming languages.
\end{enumerate}

\section*{Assessments} 
\renewcommand{\arraystretch}{1.5}
\begin{tabular}{|c|c|c|c|}
	\hline 
	\textbf{Assessment}                                 & \textbf{Weighting} & \textbf{Due Date}            & \textbf{Learning Outcome} \\ \hline
	\small Practical & \small 20\%        & \small 13-11-2024 (Wednesday at 4.59 PM)   & \small 1                   \\ \hline
	\small Project                 & \small 80\%        & \small 13-11-2024 (Wednesday at 4.59 PM) \small  & \small 1, 2                   \\ \hline
\end{tabular}

\section*{Conditions of Assessment}
You will complete this assessment during your learner-managed time. However, there will be time during class to discuss the requirements and your progress on this assessment. This assessment will need to be completed by \textbf{Wednesday, 13 November 2024} at \textbf{4.59 PM}. 

\section*{Pass Criteria}
This assessment is criterion-referenced (CRA) with a cumulative pass mark of \textbf{50\%} over all assessments in \textbf{ID608001: Intermediate Application Development Concepts}.

\section*{Authenticity}
All parts of your submitted assessment \textbf{must} be completely your work. Do your best to complete this assessment without using an \textbf{AI generative tool}. You need to demonstrate to the course lecturer that you can meet the learning outcome(s) for this assessment. \\
 
 However, if you get stuck, you can use an \textbf{AI generative tool} to help you get unstuck, permitting you to acknowledge that you have used it. In the assessment's repository \textbf{README.md} file, please include what prompt(s) you provided to the \textbf{AI generative tool} and how you used the response(s) to help you with your work. It also applies to code snippets retrieved from \textbf{StackOverflow} and \textbf{GitHub}. \\
 
 Failure to do this may result in a mark of \textbf{zero} for this assessment.

\section*{Policy on Submissions, Extensions, Resubmissions and Resits}
The school's process concerning submissions, extensions, resubmissions and resits complies with \textbf{Otago Polytechnic | Te Pūkenga} policies. Learners can view policies on the \textbf{Otago Polytechnic | Te Pūkenga} website located at \href{https://www.op.ac.nz/about-us/governance-and-management/policies}{https://www.op.ac.nz/about-us/governance-and-management/policies}.

\section*{Submission}
You \textbf{must} submit all application files via \textbf{GitHub Classroom}. Here is the URL to the repository you will use for your submission – \href{https://classroom.github.com/a/eGYqq7-g}{https://classroom.github.com/a/eGYqq7-g}. If you do not have not one, create a \textbf{.gitignore} and add the ignored files in this resource - \href{https://raw.githubusercontent.com/github/gitignore/main/Node.gitignore}{https://raw.githubusercontent.com/github/gitignore/main/Node.gitignore}. Create a branch called \textbf{project}. The latest application files in the \textbf{project} branch will be used to mark against the \textbf{Functionality} criterion. Please test before you submit. Partial marks \textbf{will not} be given for incomplete functionality. Late submissions will incur a \textbf{10\% penalty per day}, rolling over at \textbf{5:00 PM}.\\

\section*{Extensions}
Familiarise yourself with the assessment due date. Extensions will \textbf{only} be granted if you are unable to complete the assessment by the due date because of \textbf{unforeseen circumstances outside your control}. The length of the extension granted will depend on the circumstances and \textbf{must} be negotiated with the course lecturer before the assessment due date. A medical certificate or support letter may be needed. Extensions will not be granted on the due date and for poor time management or pressure of other assessments.

\section*{Resits}
Resits and reassessments \textbf{are not} applicable in \textbf{ID608001: Intermediate Application Development Concepts}.

\section*{Instructions}

\subsection*{Functionality - Learning Outcomes 1 and 2 (60\%)}

\textbf{Milestone One - Due Wednesday, 11 September 2024 (Week 8) at 4.59 PM. \\ Code Defence - Thursday, 12 September 2024 (Week 8) at 10.00 AM.}
\begin{itemize} 
	\item \textbf{Movie Listings Application (10\%):}
	\begin{itemize}
		\item Create a new directory called \textbf{movie-listings-application} and create a new \textbf{React} application using \textbf{Vite}.
		\item Can run in development without modification.
		\item Use \textbf{React Query} and the \textbf{Fetch API} to fetch data from the \textbf{The Movie DB API}.
		\item Navigation bar using \textbf{Tailwind CSS} and \textbf{Shadcn UI} with the following movie type options:
		\begin{itemize}
			\item Trending. This is the default option.
			\item Top Rated
			\item Action
			\item Animation
			\item Comedy
		\end{itemize}
		\textbf{Note:} You will use \textbf{React Router} to navigate between the options.
		\item The endpoints for the movie types are as follows:
	\begin{itemize}
		\item Trending - \href{https://api.themoviedb.org/3/trending/all/week?api\_key=$<$API KEY$>$\&language=en-US}{https://api.themoviedb.org/3/trending/all/week?api\_key=$<$API KEY$>$\&language=en-US}
		\item Top Rated - \href{https://api.themoviedb.org/3/movie/top\_rated?api\_key=$<$API KEY$>$\&language=en-US}{https://api.themoviedb.org/3/movie/top\_rated?api\_key=$<$API KEY$>$\&language=en-US}
		\item Action - \href{https://api.themoviedb.org/3/discover/movie?api\_key=$<$API KEY$>$\&with\_genres=28}{https://api.themoviedb.org/3/discover/movie?api\_key=$<$API KEY$>$\&with\_genres=28}
		\item Animation - \href{https://api.themoviedb.org/3/discover/movie?api\_key=$<$API KEY$>$\&with\_genres=16}{https://api.themoviedb.org/3/discover/movie?api\_key=$<$API KEY$>$\&with\_genres=16}
		\item Comedy - \href{https://api.themoviedb.org/3/discover/movie?api\_key=$<$API KEY$>$\&with\_genres=35}{https://api.themoviedb.org/3/discover/movie?api\_key=$<$API KEY$>$\&with\_genres=35}
	\end{itemize}
	\textbf{Note:} Replace $<$API KEY$>$ with your \textbf{The Movie DB API} key. More information can be \\found at: \href{https://developer.themoviedb.org/docs/getting-started}{https://developer.themoviedb.org/docs/getting-started}.
	\item When an option is selected, display the title, overview, poster path and release date of the first 10 movies.
	\begin{itemize}
		\item Style each story to look like a card using \textbf{Tailwind CSS} and \textbf{Shadcn UI}.
		\item Display five stories per row.
	\end{itemize}
	\end{itemize}
\end{itemize}
	\textbf{Milestone Two - Due Wednesday, 16 October 2024 (Week 11) at 4.59 PM. \\ Code Defence - Thursday, 17 October 2024 (Week 11) at 10.00 AM.}
	\begin{itemize}
	\item \textbf{Hacker News Application (20\%):}
	\begin{itemize}
		\item Create a new directory called \textbf{hacker-news-application} and create a new \textbf{React} application using \textbf{Vite}.
		\item Can run in development without modification.
		\item Use \textbf{React Query} and \textbf{GraphQL} to fetch data from the \textbf{Hacker News API}.
		\item Navigation bar using \textbf{Tailwind CSS} and \textbf{Shadcn UI} with the following story options:
		\begin{itemize}
			\item Ask Stories. This is the default option.
			\item Best Stories
			\item Job Stories
			\item New Stories
			\item Show Stories
			\item Top Stories
		\end{itemize}
		\textbf{Note:} You will use \textbf{React Router} to navigate between the options.
	\item The endpoints for the stories are as follows:
	\begin{itemize}
		\item Ask Stories - \href{https://hacker-news.firebaseio.com/v0/askstories.json?print=pretty}{https://hacker-news.firebaseio.com/v0/askstories.json?print=pretty}
		\item Best Stories - \href{https://hacker-news.firebaseio.com/v0/beststories.json?print=pretty}{https://hacker-news.firebaseio.com/v0/beststories.json?print=pretty}
		\item Job Stories - \href{https://hacker-news.firebaseio.com/v0/jobstories.json?print=pretty}{https://hacker-news.firebaseio.com/v0/jobstories.json?print=pretty}
		\item New Stories - \href{https://hacker-news.firebaseio.com/v0/newstories.json?print=pretty}{https://hacker-news.firebaseio.com/v0/newstories.json?print=pretty}
		\item Show Stories - \href{https://hacker-news.firebaseio.com/v0/showstories.json?print=pretty}{https://hacker-news.firebaseio.com/v0/showstories.json?print=pretty}
		\item Top Stories - \href{https://hacker-news.firebaseio.com/v0/topstories.json?print=pretty}{https://hacker-news.firebaseio.com/v0/topstories.json?print=pretty}
	\end{itemize}
	\item When an option is selected, display the title of the first 25 stories.
	\begin{itemize}
		\item Style each story to look like a card using \textbf{Tailwind CSS} and \textbf{Shadcn UI}.
		\item Display five stories per row.
	\end{itemize}
	\item When a story is clicked, the user will be navigate to a page that displays the following information:
	\begin{itemize}
		\item By
		\item Kids. \textbf{Note:} This is an array of ids. If the array is empty, display \textbf{N/A}. Display the first five ids as URLs in this format: \href{https://hacker-news.firebaseio.com/v0/item/$<$Id$>$.json?print=pretty}{https://hacker-news.firebaseio.com/v0/item/$<$Id$>$.json?print=pretty}. When clicked, these URLs will open in a new tab. 
		\item Score
		\item Time. \textbf{Note:} Convert the time to a readable format.
		\item Title
		\item Type 
		\item URL. \textbf{Note:} When clicked, this URL will open in a new tab.
	\end{itemize}
	This information is fetched from the following endpoint:\\
	\href{https://hacker-news.firebaseio.com/v0/item/$<$Id$>$.json?print=pretty}{https://hacker-news.firebaseio.com/v0/item/$<$Id$>$.json?print=pretty}.
	\item On a page, display the top 15 leaders. You can get this information from the following URL:\\
	\href{https://news.ycombinator.com/leaders}{https://news.ycombinator.com/leaders}. \textbf{Note:} You need to manually retrieve the information from the URL. 
	\item When a leader is clicked, the user will be navigate to a page that displays the following information:
	\begin{itemize}
		\item About
		\item Created. \textbf{Note:} Convert the time to a readable format.
		\item Id
		\item Karma
		\item Submitted. \textbf{Note:} This is an array of ids. Display the first 5 ids as URLs in this format:\\
		 \href{https://hacker-news.firebaseio.com/v0/item/$<$Id$>$.json?print=pretty}{https://hacker-news.firebaseio.com/v0/item/$<$Id$>$.json?print=pretty}.
	\end{itemize}
	This information is fetched from the following endpoint:\\
	\href{https://hacker-news.firebaseio.com/v0/user/$<$Id$>$.json?print=pretty}{https://hacker-news.firebaseio.com/v0/user/$<$Id$>$.json?print=pretty}.
	\end{itemize}
\end{itemize}
	\textbf{Milestone Three - Due Wednesday, 18 November 2024 (Week 15) at 4.59 PM. \\ Code Defence - Thursday, 19 November 2024 (Week 15) at 10.00 AM.}
	\begin{itemize}
	\item \textbf{Trivia Application (20\%):}
	\begin{itemize}
		\item Create a new directory called \textbf{trivia-application} and create a new \textbf{Next.js} application.
		\item Can run in development without modification.
		\item Use \textbf{React Query} and \textbf{GraphQL} to fetch data from the \textbf{OpenTDB API}.
		\item Create a form with the following inputs:
		\begin{itemize}
			\item Amount. This is the number of questions to retrieve. The default value is 10.
			\item Category. This is a dropdown list. The options can be retrieved from the following endpoint:\\
			\href{https://opentdb.com/api\_category.php}{https://opentdb.com/api\_category.php}. The default option is \textbf{Any Category}.			
			\item Difficulty. This is a dropdown list with the following options:
			\begin{itemize}
				\item Any Difficulty. This is the default option.
				\item Easy. 
				\item Medium
				\item Hard
			\end{itemize}
			\item Type. This is a dropdown list with the following options:
			\begin{itemize}
				\item Any Type. This is the default option.
				\item Multiple Choice
				\item True/False
			\end{itemize}
		\end{itemize}
		\item When the form is submitted, perform a request to the \href{OpenTDB API}{https://opentdb.com/api\_config.php}. 
		\item Use the response to display the questions and answers in a quiz format. 
		\item Create a form that allows the user to submit their answers. After the user submits their answers, display their score and the correct answers.
	\end{itemize} 
\end{itemize}

\subsection*{Code Quality and Best Practices - Learning Outcome 1 (45\%)}
\begin{itemize}
	\item A \textbf{Node.js} \textbf{.gitignore} file is used.
	\item Environment variables' key is stored in the \textbf{.env.example} file. 
  	\item Appropriate naming of files, variables, functions and components.
	\item Idiomatic use of \textbf{TypeScript}, control flow, data structures and in-built functions.
	\item Efficient algorithmic approach.
	\item Sufficient modularity.
	\item Each file has a \textbf{JSDoc} header comment located at the top of the file.
	\item Code is formatted.
	\item No dead or unused code. 
\end{itemize}

\subsection*{Documentation and Git Usage - Learning Outcomes 1 (5\%)}
\begin{itemize}
	\item A \textbf{GitHub} project board or issues to help you organise and prioritise your development work. The course lecturer needs to see consistent use of the \textbf{GitHub} project board or issues for the duration of the assessment.
	\item Provide the following in your repository \textbf{README.md} file:
	\begin{itemize}
		\item How do you setup the environments, i.e., after the repository is cloned?
		\item How do you run your \textbf{frontend applications} locally?
		\item How do you check your code?
		\item How do you format your code?
	\end{itemize}
    \item Use of \textbf{Markdown}, i.e., headings, bold text, code blocks, etc.
    \item Correct spelling and grammar.
    \item Your \textbf{Git commit messages} should:
    \begin{itemize}
      \item Reflect the context of each functional requirement change.
      \item Be formatted using an appropriate naming convention style.
    \end{itemize}
\end{itemize}

\subsection*{Additional Information}
\begin{itemize}
    \item If you do not attend a milestone code defence, you will receive a mark of \textbf{zero} for that milestone.
	\item \textbf{Do not} rewrite your \textbf{Git} history. It is important that the course lecturer can see how you worked on your assessment over time. 
    \item You need to show the course lecturer the initial \textbf{GitHub} project board or issues before you start your development work. Following this, you need to show the course lecturer your \textbf{GitHub} project board or issues at the end of each week.
\end{itemize} 
\end{document}
\end{document}
